\section*{Chapitre \ref{ch:s2:dispersion-and-correlation} --- Dispersion et corrélation}

\begin{solution}{exo:s2:dispersion-and-correlation:1}
Pour de nombreux membres à pondération égale, le terme de variance propre s'annule et $\bar\rho \approx \sigma_I^2/\sigma^2 = 0.04/0.09 = 0.444$.
\end{solution}

\begin{solution}{exo:s2:dispersion-and-correlation:2}
Par $\sigma_I/\sum_i w_i\sigma_i = 0.2/0.3 = 0.667$.
\end{solution}

\begin{solution}{exo:s2:dispersion-and-correlation:3}
Sur les pondérations de l'indice, le portefeuille est long de plus de variance des membres qu'il n'est court de variance de l'indice~; il gagne donc lorsque la volatilité monte de façon générale (krachs) et paie
la prime des membres pendant les mois calmes~: de petites pertes régulières, des gains rares.
\end{solution}

\begin{solution}{exo:s2:dispersion-and-correlation:4}
$0.12 - 0.83 = -0.71$ par unité. Le prix d'exercice reflétait les anticipations d'un mois calme~: la corrélation implicite est faible lorsque les marchés sont calmes, et c'est précisément à ce moment que
débutent les krachs qui font monter la corrélation.
\end{solution}

\begin{solution}{exo:s2:dispersion-and-correlation:5}
La prime est mesurée sur l'indice et sur l'ensemble des membres~; la négocier suppose des options sur chaque membre, dont les fourchettes et les coûts de couverture sont importants par rapport à
la prime~: une limite à l'arbitrage, ce qui explique aussi que la prime puisse persister.
\end{solution}

\begin{solution}{exo:s2:dispersion-and-correlation:6}
Les options sur indice sont bon marché par rapport aux options sur actions individuelles~: on anticipe que les membres s'écarteront les uns des autres. Vendre la corrélation à partir de ce niveau laisse peu de prime et
expose à une forte perte si les membres commencent à évoluer ensemble, comme ils le font lors d'une chute des cours.
\end{solution}

\begin{solution}{exo:s2:dispersion-and-correlation:7}
La prime des membres étant égale à celle de l'indice, la corrélation implicite moyenne vaut 0.258 contre une \omterm{def:s2:dispersion-and-correlation:realised}{corrélation réalisée} de 0.252~: la prime de corrélation disparaît presque
(0.006). Le ratio de Sharpe du swap de corrélation tombe à 0.17, le portefeuille neutre en vega perd ($-1.84$), et le portefeuille sur les pondérations de l'indice perd à chaque mois
($-17.3$), car il paie la prime de variance complète sur davantage de variance des membres qu'il n'en vend sur l'indice.
\end{solution}

\begin{solution}{exo:s2:dispersion-and-correlation:8}
Gagner dans 77~\% des mois est la forme typique d'un trade de volatilité courte~: les pertes surviennent lors des autres mois, et le premier krach a coûté au swap de corrélation six
écarts-types mensuels. Un effet de levier choisi d'après la fréquence des gains ne résisterait pas à l'ampleur des pertes~; il faut le dimensionner par un scénario de stress sur un mois de krach.
\end{solution}

\begin{solution}{pb:s2:dispersion-and-correlation:1}
\begin{enumerate}
\item $\bar\rho = (\sigma_I^2 - \sum w_i^2\sigma_i^2)/\sum_{i \ne j} w_i w_j\sigma_i\sigma_j$~: la corrélation unique qui réconcilie les
volatilités de l'indice et celles des membres.
\item La même formule appliquée aux variances réalisées~; corrélation implicite moins \omterm{def:s2:dispersion-and-correlation:realised}{corrélation réalisée}.
\item À partir de la volatilité implicite du SPX et des volatilités implicites des 50 plus gros membres, avec la décomposition de Markowitz.
\item La corrélation est un risque qui est tarifé et qui explique les rendements des options~; une stratégie qui la vend a un alpha élevé sans frictions et n'en a aucun avec frictions.
\item 37.0 en moyenne, 10.8 à 69.7 pour les 90~\% centraux, 96.6 au sommet (24 octobre 2008), 70.4 et 72.5 en octobre 2008 et en mars 2020.
\item Indice de dispersion du S\&P 500 de Cboe~: 25.5 en moyenne depuis 2014, 58.9 au sommet le 18 mars 2020.
\item Il reste peu de prime et la perte si la corrélation monte est importante.
\item Une \omterm{def:s2:dispersion-and-correlation:realised}{corrélation réalisée} lissée est dans la plupart des cas légèrement inférieure à l'implicite.
\item Les notionnels des membres sur les pondérations de l'indice, ou mis à l'échelle pour égaler le vega de l'indice.
\item $-0.02$, 0.65 et 1.74~; asymétrie 2.67, $-5.87$ et $-1.67$.
\item Il paie la prime des membres sur davantage de variance qu'il n'en vend.
\item Il paie directement la corrélation implicite moins la \omterm{def:s2:dispersion-and-correlation:realised}{corrélation réalisée}, sans exposition au niveau de variance.
\item \textbf{Résultat nommé.} Le portefeuille de dispersion neutre en vega gagne la prime de corrélation (implicite 0.311 contre réalisée 0.252) avec un ratio de Sharpe de
0.65, et perd 11.3 écarts-types mensuels lorsque la corrélation passe de 0.12 à 0.83 lors du premier krach~; le swap de corrélation gagne 1.74 et perd
0.71 par unité.
\item De 0.12 à 0.83, de 0.12 à 0.57, de 0.73 à 0.84.
\item La corrélation implicite était déjà élevée lorsqu'il a commencé.
\item Par la perte dans un mois de krach, la corrélation tendant vers un.
\item Coûts des options sur actions individuelles, mois de krach, \omterm{def:s2:dispersion-and-correlation:realised}{corrélation réalisée} telle que le contrat la définit.
\item Dispersion pondérée par le vega.
\item Le chapitre~1 vendait de la volatilité~; la dispersion vend la part de la volatilité de l'indice qui provient de la corrélation.
\item Assurance contre le fait que les membres évoluent ensemble.
\end{enumerate}
\end{solution}

\begin{solution}{iq:s2:dispersion-and-correlation:1}
La corrélation moyenne unique qui rend la variance implicite de l'indice égale à la somme pondérée des variances implicites et des covariances des membres~:
$(\sigma_I^2 - \sum w_i^2\sigma_i^2)/\sum_{i \ne j} w_i w_j\sigma_i\sigma_j$ avec les volatilités implicites.
\end{solution}

\begin{solution}{iq:s2:dispersion-and-correlation:2}
Vendre des options ou des variances sur l'indice et acheter celles des membres, pondérées par le vega, le theta ou les pondérations de l'indice~; la stratégie rapporte la prime de corrélation et
perd lorsque les membres évoluent ensemble~; elle comporte aussi un risque d'événement propre à chaque action ainsi que la liquidité des membres et les coûts de couverture.
\end{solution}

\begin{solution}{iq:s2:dispersion-and-correlation:3}
Les options sur indice assurent contre les baisses générales, lorsque la corrélation monte et que la diversification échoue~; les investisseurs paient cette assurance, et l'offre d'options sur indice
(overwriting) est plus faible que la demande, alors que les options sur actions individuelles sont souvent offertes par des vendeurs d'options couvertes.
\end{solution}

\begin{solution}{iq:s2:dispersion-and-correlation:4}
Appliquer des mois de krach historiques (2008, 2020) et un mouvement hypothétique de la corrélation implicite et réalisée vers un, avec des volatilités des membres à la hausse~; il faut
comparer la perte aux limites et à la liquidité du dénouement des jambes sur actions individuelles.
\end{solution}

\begin{solution}{iq:s2:dispersion-and-correlation:5}
Les surfaces de volatilité des options sur l'indice et sur les membres, les pondérations de l'indice et les changements de composition, les événements d'entreprise (résultats, fusions) par membre, et les emprunts de titres pour les couvertures~; des cours
affichés obsolètes sur actions individuelles, et des erreurs de composition faussent la corrélation implicite.
\end{solution}

\begin{solution}{iq:s2:dispersion-and-correlation:6}
Avec $w_i = 1/N$ et un $\sigma$ commun, $\sigma_I^2 = \sum_i\sigma^2/N^2 + \sum_{i \ne j}\rho\sigma^2/N^2 = \sigma^2/N + \rho\sigma^2(N - 1)/N = \sigma^2(\rho +
(1 - \rho)/N)$. Pour $N$ grand, $\rho \approx \sigma_I^2/\sigma^2$.
\end{solution}
